% 05_architecture.tex
\section{Architecture}

% Required TikZ libraries for fit/background scopes used in this chapter's diagrams.
\usetikzlibrary{fit,backgrounds}

\subsection{Architectural Overview: MVC Intuition}

The application strictly adheres to the Model-View-Controller (MVC) design pattern. This architectural choice separates the internal representations of data structures, the logic that manipulates them, and the graphical interfaces presented to the user. This separation of concerns allows for a highly modular, maintainable, and extensible codebase.

The architecture is divided into three interconnected categories of components:
\begin{itemize}
    \item \textbf{Controller}: Orchestrates the application. The central \texttt{AppEngine} class serves as the primary coordinator, processing user inputs and coordinating changes across the system. It delegates playback controls to the \texttt{Playback} controller, which manages the flow of time for step-by-step animations.
    \item \textbf{Model}: Represents the mathematical data structures, their algorithms, and their visual states over time. The \texttt{IVisualizable} interface acts as the base abstraction for all concrete implementations (such as Linked Lists, Graphs, and Heaps). When an algorithm is executed, the structure generates a \texttt{Timeline} composed of discrete \texttt{Frame}s. Each frame captures the exact state of the structure at a given execution step.
    \item \textbf{View}: Responsible for rendering the visual representation of the \texttt{Model}'s state. The \texttt{UIManager} handles modern graphical user interface elements via Dear ImGui, providing interactive menus and code panels. The \texttt{Renderer} utilizes SFML to draw graphical shapes based on the data contained in the current \texttt{Frame}.
\end{itemize}

\begin{figure}[h!]
    \centering
    \resizebox{\textwidth}{!}{
    \begin{tikzpicture}[
        node distance=1.5cm and 2cm,
        box/.style={rectangle, draw, thick, minimum width=3.5cm, minimum height=1cm, align=center, fill=white, rounded corners=2pt, font=\small},
        groupbox/.style={rectangle, draw=gray, dashed, thick, inner sep=15pt},
        arrow/.style={-{Stealth[scale=1.2]}, thick, draw=black!80},
        dashed_arrow/.style={-{Stealth[scale=1.2]}, dashed, thick, draw=black!80},
        labelnode/.style={font=\bfseries}
    ]
    
    % ================== CONTROLLER ==================
    \node (appEngine) [box, fill=green!10] {\textbf{AppEngine}\\(Coordinator)};
    \node (dataManager) [box, fill=green!10, left=of appEngine, xshift=-0.5cm] {DataManager};
    \node (playback) [box, fill=green!10, right=of appEngine, xshift=0.5cm] {Playback};
    
    \begin{scope}[on background layer]
        \node (controllerGroup) [groupbox, fill=green!5, fit=(appEngine) (dataManager) (playback)] {};
        \node [above left=0pt of controllerGroup.north west, labelnode, color=green!60!black] {Controller Layer};
    \end{scope}
    
    % ================== VIEW ==================
    \node (uimanager) [box, fill=red!10, below left=3.5cm and -1.5cm of appEngine] {\textbf{UIManager}\\(ImGui)};
    \node (renderer) [box, fill=red!10, below=of uimanager] {\textbf{Renderer}\\(SFML)};
    
    \begin{scope}[on background layer]
        \node (viewGroup) [groupbox, fill=red!3, fit=(uimanager) (renderer)] {};
        \node [above left=0pt of viewGroup.north west, labelnode, color=red!70!black] {View Layer};
    \end{scope}
    
    % ================== MODEL ==================
    \node (ivisualizable) [box, fill=blue!10, below right=3.5cm and -0.5cm of appEngine] {\textbf{IVisualizable}\\(Base Interface)};
    \node (structures) [box, fill=blue!5, below=of ivisualizable] {Data Structures\\(List, Graph, Heap, etc.)};
    
    \node (timeline) [box, fill=cyan!10, right=of ivisualizable, xshift=1cm] {Timeline \& Frames};
    \node (payloads) [box, fill=cyan!5, below=of timeline] {IPayload};
    \node (ipayloadvis) [box, fill=cyan!5, right=of payloads, xshift=0.5cm] {IPayloadVisitor};
    
    \begin{scope}[on background layer]
        \node (modelGroup) [groupbox, fill=blue!3, fit=(ivisualizable) (structures) (timeline) (payloads) (ipayloadvis)] {};
        \node [above left=0pt of modelGroup.north west, labelnode, color=blue!70!black] {Model Layer};
    \end{scope}
    
    % ================== UTILITIES ==================
    \node (utilities) [box, fill=yellow!10, below=1.5cm of controllerGroup] {Utilities\\(Pseudocode, config)};

    % ================== RELATIONSHIPS ==================
    % Model internal
    \draw [arrow] (structures) -- (ivisualizable);
    \draw [dashed_arrow] (ivisualizable) -- node[above, font=\scriptsize] {generates} (timeline);
    \draw [arrow] (timeline) -- node[right, font=\scriptsize] {owns} (payloads);
    \draw [dashed_arrow] (ipayloadvis) -- node[above, font=\scriptsize] {visits} (payloads);
    
    % Controller connections
    \draw [arrow] (appEngine) -- (dataManager);
    \draw [arrow] (appEngine) -- (playback);
    \draw [arrow] (appEngine) -- node[above right, font=\scriptsize, sloped, pos=0.6] {routes actions} (modelGroup.north west);
    \draw [dashed_arrow] (playback.east) -| ++(1.5,0) |- node[above, font=\scriptsize, near end] {traverses} (timeline.east);
    
    % View connections
    \draw [arrow] (appEngine) -- node[above left, font=\scriptsize, sloped, pos=0.6] {updates} (viewGroup.north east);
    
    % Cross-layer interactions
    \draw [arrow] (renderer.south) |- ++(0,-1cm) -| node[right, font=\scriptsize, near end] {implements} (ipayloadvis.south);
    \draw [dashed_arrow] (renderer.east) -- node[above, font=\scriptsize] {renders} (payloads.west);
    
    \draw [dashed_arrow] (uimanager.south) |- ++(0,-2cm) -| node[right, font=\scriptsize, near end] {queries state} (timeline.south);
    
    \draw [arrow] (uimanager.north) |- node[above, font=\scriptsize, near end] {reads context} (utilities.west);
    \end{tikzpicture}
    }
    \caption{Application Architecture Diagram detailing the relationships between MVC components.}
    \label{fig:arch_diagram}
\end{figure}

\subsection{Detailed System Design and Interfaces}

A cornerstone of the application's design is its extensive use of abstract interfaces. This approach enforces loose coupling, allowing the front-end rendering logic to remain agnostic to the internal implementations of the back-end data structures.

\subsubsection{The \texttt{IVisualizable} Abstraction}
Every data structure in the application inherits from the \texttt{IVisualizable} interface. This ensures that the \texttt{AppEngine} can interact with an AVL Tree, a Graph, or a Singly Linked List using a unified API. When a user action is dispatched (e.g., inserting a node), the specific data structure executes its algorithm and generates a \texttt{Timeline}.

\subsubsection{Timeline and Frames}
Instead of immediately drawing the results of an algorithm, the models generate discrete states.
\begin{itemize}
    \item \textbf{Timeline}: A collection of \texttt{Frame} objects, representing the complete history of an algorithm's execution from start to finish.
    \item \textbf{Frame}: Captures a static snapshot of the visualization at a specific moment in time. A \texttt{Frame} holds the highlighted pseudocode line, variable states, and a list of graphical \texttt{IPayload} objects that need to be drawn.
\end{itemize}
This design effectively isolates the algorithmic logic from the real-time rendering loop. It natively supports playback features such as pausing, rewinding, and fast-forwarding, simply by selecting a different \texttt{Frame} within the \texttt{Timeline}.

\begin{figure}[h!]
    \centering
    \begin{tikzpicture}[
        node distance=1.5cm and 2cm,
        base/.style={ellipse, draw, thick, minimum width=2.5cm, minimum height=1cm, align=center},
        process/.style={rectangle, draw, thick, minimum width=3cm, minimum height=1cm, align=center, fill=blue!5},
        decision/.style={diamond, aspect=1.5, draw, thick, minimum width=2.5cm, minimum height=1cm, align=center, fill=orange!5},
        arrow/.style={-{Stealth[scale=1.2]}, thick}
    ]
    
    % Nodes
    \node (start) [base, fill=gray!10] {Main Loop};
    \node (poll) [process, below=of start] {Process Input\\(SFML Events)};
    \node (update) [process, below=of poll] {Update System State\\(UI state, Playback)};
    \node (handleAction) [decision, below=of update, yshift=-1cm] {Action\\Requested?};
    
    \node (genTimeline) [process, right=of handleAction, xshift=1cm] {Model executes action\\Generates Timeline};
    \node (assignTimeline) [process, above=of genTimeline] {Assign Timeline\\to Playback};
    
    \node (render) [process, below=of handleAction, yshift=-1cm] {Render Current Frame\\(Renderer \& UI)};
    
    % Arrows
    \draw [arrow] (start) -- (poll);
    \draw [arrow] (poll) -- (update);
    \draw [arrow] (update) -- (handleAction);
    
    \draw [arrow] (handleAction) -- node[anchor=north] {Yes (Insert/Delete/\dots)} (genTimeline);
    \draw [arrow] (handleAction) -- node[anchor=east] {No} (render);
    
    \draw [arrow] (genTimeline) -- (assignTimeline);
    \draw [arrow] (assignTimeline) |- (render);
    
    \draw [arrow] (render) -- ++(-4,0) |- (poll);
    
    \end{tikzpicture}
    \caption{Application Execution Flow Loop.}
    \label{fig:app_flow}
\end{figure}

\subsection{Visitor Design Pattern for Rendering}

A major challenge in a visualization system is handling the heterogeneous types of data that need to be drawn. A Singly Linked List requires rendering distinct nodes and directional arrows, while an Adjacency Matrix necessitates rendering a grid. If the \texttt{Renderer} were to use a long series of \texttt{if-else} statements to determine the type of an object before drawing it, the code would quickly become unmaintainable. Furthermore, giving the \texttt{Model} the responsibility to render itself would violate the MVC principles.

To solve this problem cleanly, the application employs the \textbf{Visitor Design Pattern}. 

\subsubsection{Pattern Implementation}
The pattern relies on two primary abstractions:
\begin{enumerate}
    \item \textbf{\texttt{IPayloadVisitor}}: An interface declaring a \texttt{visit()} method for every possible concrete payload type (e.g., \texttt{visit(const LinkedListPayload\&)}, \texttt{visit(const GraphPayload\&)}). The \texttt{Renderer} implements this interface.
    \item \textbf{\texttt{IPayload}}: An interface representing a drawing instruction generated by the Model. It requires an \texttt{accept(IPayloadVisitor\& visitor)} method.
\end{enumerate}

When a \texttt{Frame} is rendered, the \texttt{Renderer} iterates through the \texttt{Frame}'s payload list and calls \texttt{payload->accept(*this)}. Through dynamic dispatch, the concrete payload's \texttt{accept} method is invoked. This method then immediately calls \texttt{visitor.visit(*this)}, effectively passing its strongly-typed self back to the \texttt{Renderer}.

\begin{figure}[h!]
    \centering
    \resizebox{\textwidth}{!}{
    \begin{tikzpicture}[
        node distance=2cm and 2cm,
        box/.style={rectangle, draw, thick, minimum width=3cm, minimum height=1cm, align=center, fill=white, rounded corners=2pt},
        interface/.style={box, fill=blue!5, font=\ttfamily},
        class/.style={box, fill=green!5, font=\ttfamily},
        arrow/.style={-{Stealth[scale=1.2]}, thick},
        dashed_arrow/.style={-{Stealth[scale=1.2]}, dashed, thick}
    ]
    
    % Nodes
    \node (iv) [interface] {IPayloadVisitor \\ \rule{2.6cm}{0.4pt} \\ + visit(LinkedListPayload) \\ + visit(GraphPayload) \\ \dots};
    \node (renderer) [class, below=of iv, yshift=-0.5cm] {Renderer \\ \rule{2.6cm}{0.4pt} \\ + visit(LinkedListPayload) \\ + visit(GraphPayload) \\ \dots};
    
    \node (ip) [interface, right=of iv, xshift=5cm] {IPayload \\ \rule{2.6cm}{0.4pt} \\ + accept(IPayloadVisitor) \\ + clone()};
    \node (llp) [class, below=of ip, xshift=-2.5cm, yshift=-0.5cm] {LinkedListPayload \\ \rule{2.6cm}{0.4pt} \\ + accept(v: IPayloadVisitor) \\ \textit{\{ v.visit(*this) \}}};
    \node (gp) [class, below=of ip, xshift=2.5cm, yshift=-0.5cm] {GraphPayload \\ \rule{2.6cm}{0.4pt} \\ + accept(v: IPayloadVisitor) \\ \textit{\{ v.visit(*this) \}}};
    
    % Inheritance arrows
    \draw [{Stealth[scale=1.2, fill=white]}-, thick] (iv) -- (renderer);
    \draw [{Stealth[scale=1.2, fill=white]}-, thick] (ip) -- (llp);
    \draw [{Stealth[scale=1.2, fill=white]}-, thick] (ip) -- (gp);
    
    % Dependency arrows
    \draw [dashed_arrow] (llp.west) -- node[above, font=\scriptsize] {calls} (renderer.east);
    \draw [dashed_arrow] (gp.south) |- ++(0,-0.8cm) -| node[above, near start, font=\scriptsize] {calls} (renderer.south);
    
    \end{tikzpicture}
    }
    \caption{UML Class Diagram illustrating the Visitor Pattern for Payload rendering.}
    \label{fig:visitor_pattern}
\end{figure}

This architecture guarantees that the \texttt{Model} never references SFML or rendering specifics, while the \texttt{Renderer} avoids costly runtime type checks (\texttt{dynamic\_cast}), resulting in safe, robust, and highly decoupled code.
